\chapter{Valutazione dell'esperienza di stage}
\label{chap:valutazione}

Questo capitolo effettua una valutazione retrospettiva del percorso di stage, sul piano del raggiungimento degli obiettivi, della crescita professionale personale, e del confronto tra le competenze richieste dallo stage e quelle fornite dal corso di studi.

\section{Raggiungimento degli obiettivi}
\label{sec:raggiungimento-obiettivi}

\todoplaceholder{grado di soddisfacimento degli obiettivi, su base fattuale.}

% TODO: valutazione, basata su dati di fatto e non su opinioni personali, del
% grado di soddisfacimento degli obiettivi indicati nel capitolo~\ref{chap:motivazioni-obiettivi},
% sia dal punto di vista dello studente sia da quello dell'organizzazione ospitante.


\section{Crescita professionale}
\label{sec:crescita-professionale}

\todoplaceholder{conoscenze e abilità acquisite durante lo stage.}

% TODO: riflessione su cosa lo stage ha effettivamente portato in termini di
% conoscenze e abilità acquisite, con riferimento concreto alle attività
% descritte nel capitolo~\ref{chap:sviluppo-estensione}.


\section{Confronto tra competenze richieste e percorso di studi}
\label{sec:confronto-competenze}

\todoplaceholder{distanza tra competenze richieste e percorso di studi.}

% TODO: distanza tra le competenze richieste a inizio stage per svolgere il
% lavoro e quelle fornite dal corso di studi, con eventuale segnalazione di
% lacune riscontrate.
